% =============================================================
%  CS102-F26 Digital Design and Computer Architecture I
%  Lecture 10 -- Finite State Machines (Unit 10)
%  Follows Chapter 5 (introduction) and Section 5.1 of
%  A Student's Computer Architecture Book I.
%
%  Build the slides and the print handout:  npm run materials -- lecture-10
% =============================================================
\makeatletter
\def\input@path{{templates/latex/}{../}}
\makeatother

\documentclass[10pt]{wosm-slides}
\usepackage{cs102-lectures}[2026/10/20]

\title[Lecture 10 \textperiodcentered\ Finite State Machines]{Finite State Machines}
\subtitle{Lecture 10: sequential logic, Moore and Mealy machines, and simplification}

% ---------- Helpers for this lecture --------------------------
\tikzset{
  fsm/.style={st, minimum size=1.35cm, font=\footnotesize, align=center},
  initial arrow/.style={arrow, {Stealth[length=5pt]}-},
  every loop/.style={arrow, looseness=5},
}
% A block for block diagrams.
\tikzset{block/.style={draw=wosm-text, line width=0.6pt, minimum width=2.1cm,
  minimum height=0.9cm, font=\scriptsize, align=center}}

% A state table row: \srow{<state>}{<next on 0>}{<next on 1>}{<output>}
\newcommand*\srow[4]{#1 & #2 & #3 & #4 \\}

% =============================================================
\begin{document}

\maketitle

\begin{frame}{Today}
  \tableofcontents
\end{frame}

% =============================================================
\section{Sequential logic}

\begin{frame}{Beyond combinational logic}
  \begin{itemize}
    \item A combinational circuit computes its outputs from its \emph{present}
      inputs only (Lecture 07).
    \item A \alert{sequential} circuit may also take the \emph{past} into account.
      Combinational logic is the special case that ignores it.
    \item The sequential circuits a computer needs most: the \textbf{clock}, the
      \textbf{memory} and the \textbf{counter} \unitref{11--13}.
  \end{itemize}

  \bigskip
  \begin{remark}
    Sequential circuits come in two kinds. \emph{Synchronous} ones change state only
    at the ticks of a clock; \emph{asynchronous} ones whenever their inputs change.
    The difference becomes clear with the clock and memory in the next lectures.
  \end{remark}
\end{frame}

\begin{frame}{The turnstile}
  \begin{columns}[c, totalwidth=\textwidth]
    \column{0.5\textwidth}
      \begin{itemize}
        \item Two actions: use a ticket (\emph{pass}) and turn the arm
          (\emph{push}).
        \item Will the arm turn if someone pushes it? The present action alone
          cannot say: it depends on whether a ticket was used before.
        \item The turnstile must remember one fact about the past: it is
          \alert{locked} or \alert{unlocked}.
      \end{itemize}
    \column{0.46\textwidth}
      \centering
      \begin{tikzpicture}[node distance=2.4cm]
        \node[fsm] (l) {Locked};
        \node[fsm, right=of l] (u) {Unlocked};
        \draw[arrow] (l) to[bend left=25] node[above, edgelabel] {pass} (u);
        \draw[arrow] (u) to[bend left=25] node[below, edgelabel] {push} (l);
        \draw (l) edge[loop above] node[above, edgelabel] {push} (l);
        \draw (u) edge[loop above] node[above, edgelabel] {pass} (u);
        \draw[initial arrow] (l.west) -- ++(-0.5,0);
      \end{tikzpicture}
  \end{columns}

  \bigskip
  \begin{definition}[State]
    A \emph{state} is a stored summary of the past: everything that must be
    remembered to decide what happens next.
  \end{definition}
\end{frame}

\begin{frame}{The mystic box}
  A box with two buttons, \emph{On} and \emph{Off}. A pressed button sends 1, a
  released one 0. We never press both. The box should turn green on \emph{On}, red
  on \emph{Off}, and keep its colour otherwise.

  \medskip
  \begin{center}
    \small
    \begin{tabular}{@{}l|*{12}{c}@{}}
      Time & 0 & 1 & 2 & 3 & 4 & 5 & 6 & 7 & 8 & 9 & 10 & 11 \\
      \midrule
      Input (On, Off) & 10 & 00 & 00 & 00 & 01 & 00 & 00 & 10 & 00 & 01 & 01 & 00 \\
      Colour & G & G & G & G & R & R & R & G & G & R & R & R \\
    \end{tabular}
  \end{center}

  \medskip
  \begin{itemize}
    \item Input 00 appears at times 1, 2, 3, 5, 6, 8 and 11, with colour G at 1, 2,
      3, 8 and R at 5, 6, 11.
    \item A combinational circuit gives one output per input. Make 00 green, and
      times 5, 6, 11 are wrong; make it red, and times 1, 2, 3, 8 are wrong.
  \end{itemize}
  \wosmlabel{The box needs a state: the colour it showed last.}
\end{frame}

% =============================================================
\section{Finite state machines}

\begin{frame}{Finite state machines}
  \begin{definition}[Finite state machine]
    A \emph{finite state machine} (FSM) has a finite set of states, one of them the
    \emph{initial} state, inputs and outputs, and two rules:
    \begin{itemize}
      \item the \emph{transition} rule: next state $= f(\text{state}, \text{input})$;
      \item the \emph{output} rule, in one of two kinds:
        \begin{itemize}
          \item \textbf{Moore}: output $= g(\text{state})$;
          \item \textbf{Mealy}: output $= g(\text{state}, \text{input})$.
        \end{itemize}
    \end{itemize}
  \end{definition}

  \medskip
  The rule $s_{t+1} = f(s_t, x_t)$ from Lecture 01 was the transition rule of an FSM.
  The push-button lamp there was one, with two states.
\end{frame}

\begin{frame}{The mystic box as a Moore machine}
  \begin{columns}[c, totalwidth=\textwidth]
    \column{0.5\textwidth}
      \centering
      \begin{tikzpicture}[node distance=2.4cm]
        \node[fsm] (g) {Green\\[-1pt]/ G};
        \node[fsm, right=of g] (r) {Red\\[-1pt]/ R};
        \draw[arrow] (g) to[bend left=25] node[above, edgelabel] {01} (r);
        \draw[arrow] (r) to[bend left=25] node[below, edgelabel] {10} (g);
        \draw (g) edge[loop above] node[above, edgelabel] {00, 10} (g);
        \draw (r) edge[loop above] node[above, edgelabel] {00, 01} (r);
        \draw[initial arrow] (g.west) -- ++(-0.5,0);
      \end{tikzpicture}
    \column{0.46\textwidth}
      \begin{itemize}
        \item Each state shows its name and its output: \emph{state / output}.
        \item Arrows are labelled with the input (On, Off) that takes them.
        \item The arrow from nowhere marks the initial state.
      \end{itemize}
  \end{columns}

  \bigskip
  \begin{center}
    \small
    \begin{tabular}{@{}l|ccc|c@{}}
      \toprule
      State & next on 00 & next on 01 & next on 10 & output \\
      \midrule
      Green & Green & Red & Green & G \\
      Red & Red & Red & Green & R \\
      \bottomrule
    \end{tabular}
  \end{center}
\end{frame}

\begin{frame}{The mystic box as a Mealy machine}
  \begin{columns}[c, totalwidth=\textwidth]
    \column{0.5\textwidth}
      \centering
      \begin{tikzpicture}[node distance=2.4cm]
        \node[fsm] (g) {Green};
        \node[fsm, right=of g] (r) {Red};
        \draw[arrow] (g) to[bend left=25] node[above, edgelabel] {01 / R} (r);
        \draw[arrow] (r) to[bend left=25] node[below, edgelabel] {10 / G} (g);
        \draw (g) edge[loop above] node[above, edgelabel] {00 / G, 10 / G} (g);
        \draw (r) edge[loop above] node[above, edgelabel] {00 / R, 01 / R} (r);
        \draw[initial arrow] (g.west) -- ++(-0.5,0);
      \end{tikzpicture}
    \column{0.46\textwidth}
      \begin{itemize}
        \item The outputs move from the states to the arrows:
          \emph{input / output}.
        \item In Red, input 10 outputs G while it moves to Green.
      \end{itemize}
  \end{columns}

  \bigskip
  \begin{columns}[T, totalwidth=\textwidth]
    \column{0.47\textwidth}
      \textbf{Moore}
      \begin{itemize}
        \item Outputs change only when the state changes.
        \item No path from an input straight to an output.
      \end{itemize}
    \column{0.47\textwidth}
      \textbf{Mealy}
      \begin{itemize}
        \item Outputs react to the input at once.
        \item Often needs fewer states.
      \end{itemize}
  \end{columns}
\end{frame}

\begin{frame}{From a state machine to a circuit}
  \begin{center}
    \begin{tikzpicture}[node distance=0.9cm]
      \node[block] (n) {next-state\\logic};
      \node[block, right=of n] (s) {state\\memory};
      \node[block, right=of s] (o) {output\\logic};
      \draw[arrow] ($(n.west)+(-0.9,0.2)$) node[left, font=\scriptsize] {inputs}
        -- ($(n.west)+(0,0.2)$);
      \draw[arrow] (n) -- (s);
      \draw[arrow] (s) -- (o);
      \draw[arrow] (o.east) -- ++(0.6,0) node[right, font=\scriptsize] {outputs};
      \draw[arrow] (s.south) -- ++(0,-0.35) -| (n.south);
      \draw[arrow, dashed] ($(n.west)+(-0.45,0.2)$) |- ++(0,1.25) -| (o.north);
      \draw[arrow] ($(s.north)+(0,0.4)$) node[above, font=\scriptsize] {clock}
        -- (s.north);
    \end{tikzpicture}
  \end{center}

  \begin{itemize}
    \item The state is stored in memory, as bits \unitref{11--12}. The next-state and
      output logic are combinational.
    \item The dashed input to the output logic is there only in a Mealy machine.
    \item States become bits: with 2 states, 1 bit is enough. A one-hot code uses one
      bit per state, which makes the logic simpler (Lecture 08).
  \end{itemize}
\end{frame}

\begin{frame}{The mystic box as a circuit}
  \begin{columns}[c, totalwidth=\textwidth]
    \column{0.4\textwidth}
      \centering
      \begin{kmap}{2}{4}{Q}{\mathit{On}\,\mathit{Off}}{0,1}{00,01,11,10}
        \kmaprow{0}{0,0,X,1}\kmaprow{1}{1,0,X,1}
        \kgroup{0}{2}{1}{3}
        \kgroup[wosm-brick]{1}{-1}{1}{0}\kgroup[wosm-brick]{1}{3}{1}{4}
      \end{kmap}\par\smallskip
      \wosmlabel{next state $Q^{+}$}
    \column{0.56\textwidth}
      \begin{itemize}
        \item Encode Green as $Q = 1$, Red as $Q = 0$. The output is just $Q$.
        \item Next state: 1 after \emph{On}, 0 after \emph{Off}, otherwise $Q$. Both
          buttons at once never happens: X.
      \end{itemize}
      \[ Q^{+} = \mathit{On} + \mathit{Off}'\,Q \]
  \end{columns}

  \bigskip
  \wosmlabel{Coming back: a one-bit memory with a \enquote{set} and a \enquote{reset}
  input. This is the SR latch of the next lecture (Unit 11).}
\end{frame}

% =============================================================
\section{Simplifying state machines}

\begin{frame}{Equivalent states}
  \begin{itemize}
    \item Two states are \alert{equivalent} if no sequence of inputs can tell them
      apart: from either one, every input sequence gives the same outputs.
    \item Equivalent states can be merged into one. Fewer states need fewer bits of
      memory and simpler logic.
    \item The book works on the state table, and merges states whose rows are
      identical.
  \end{itemize}

  \bigskip
  \begin{block}{Simplifying a Moore machine}
    \begin{enumerate}
      \item Find two states $p$ and $q$ with the same output and the same next state
        for every input.
      \item They are equivalent: delete $q$, and replace $q$ by $p$ everywhere.
      \item Repeat until no such pair is left.
    \end{enumerate}
  \end{block}
\end{frame}

\begin{frame}{Example: a Moore machine}
  \begin{columns}[T, totalwidth=\textwidth]
    \column{0.32\textwidth}
      \centering
      \small
      \begin{tabular}{@{}c|cc|c@{}}
        state & 0 & 1 & out \\
        \midrule
        \srow{\alert<2>{a}}{b}{c}{1}
        \srow{b}{b}{d}{0}
        \srow{\alert<2>{c}}{b}{c}{1}
        \srow{d}{b}{c}{0}
        \srow{e}{b}{a}{1}
      \end{tabular}\par\smallskip
      \wosmlabel{the book's Table 5.3}
    \column{0.32\textwidth}
      \centering
      \small
      \uncover<2->{%
      \begin{tabular}{@{}c|cc|c@{}}
        state & 0 & 1 & out \\
        \midrule
        \srow{\alert<3>{a}}{b}{a}{1}
        \srow{b}{b}{d}{0}
        \srow{d}{b}{a}{0}
        \srow{\alert<3>{e}}{b}{a}{1}
      \end{tabular}\par\smallskip
      \wosmlabel{$c$ merged into $a$}}
    \column{0.32\textwidth}
      \centering
      \small
      \uncover<3->{%
      \begin{tabular}{@{}c|cc|c@{}}
        state & 0 & 1 & out \\
        \midrule
        \srow{a}{b}{a}{1}
        \srow{b}{b}{d}{0}
        \srow{d}{b}{a}{0}
      \end{tabular}\par\smallskip
      \wosmlabel{$e$ merged into $a$}}
  \end{columns}

  \bigskip
  \begin{itemize}
    \item<2-> $a$ and $c$: output 1, next states $b$ and $c$. Merge $c$ into $a$.
    \item<3-> Now $a$, $d$ and $e$ all go to $b$ and $a$, but only $a$ and $e$ output
      1. Merge $e$ into $a$.
    \item<3-> $b$ and $d$ have the same output, but on input 1 they go to $d$ and $a$:
      not mergeable. Five states became three.
  \end{itemize}
\end{frame}

\begin{frame}{Simplifying a Mealy machine}
  In a Mealy machine the outputs are on the transitions, so compare (input, next
  state, output) for every input.

  \medskip
  \begin{columns}[T, totalwidth=\textwidth]
    \column{0.32\textwidth}
      \centering
      \small
      \begin{tabular}{@{}c|cc@{}}
        state & 0 & 1 \\
        \midrule
        \alert<2>{a} & b / 1 & c / 1 \\
        b & b / 0 & d / 1 \\
        \alert<2>{c} & b / 1 & c / 1 \\
        \alert<3>{d} & e / 0 & c / 1 \\
        \alert<3>{e} & e / 0 & a / 1 \\
      \end{tabular}\par\smallskip
      \wosmlabel{the book's Table 5.6}
    \column{0.32\textwidth}
      \centering
      \small
      \uncover<2->{%
      \begin{tabular}{@{}c|cc@{}}
        state & 0 & 1 \\
        \midrule
        a & b / 1 & a / 1 \\
        b & b / 0 & d / 1 \\
        \alert<3>{d} & e / 0 & a / 1 \\
        \alert<3>{e} & e / 0 & a / 1 \\
      \end{tabular}\par\smallskip
      \wosmlabel{$c$ merged into $a$}}
    \column{0.32\textwidth}
      \centering
      \small
      \uncover<3->{%
      \begin{tabular}{@{}c|cc@{}}
        state & 0 & 1 \\
        \midrule
        a & b / 1 & a / 1 \\
        b & b / 0 & d / 1 \\
        d & d / 0 & a / 1 \\
      \end{tabular}\par\smallskip
      \wosmlabel{$e$ merged into $d$}}
  \end{columns}

  \bigskip
  \uncover<3->{The simplified machine gives the same outputs as the original for every
  input sequence.}
\end{frame}

\begin{frame}{When identical rows are not enough}
  \begin{columns}[c, totalwidth=\textwidth]
    \column{0.3\textwidth}
      \centering
      \small
      \begin{tabular}{@{}c|cc|c@{}}
        state & 0 & 1 & out \\
        \midrule
        \srow{p}{r}{\alert{q}}{0}
        \srow{q}{r}{\alert{p}}{0}
        \srow{r}{r}{p}{1}
      \end{tabular}
    \column{0.66\textwidth}
      \begin{itemize}
        \item The rows of $p$ and $q$ differ, yet the two states are equivalent: they
          only send each other back and forth.
        \item The full method, \emph{partition refinement}: group the states by
          output, then split any group whose states lead to different groups on
          some input. Repeat until nothing splits.
      \end{itemize}
  \end{columns}

  \bigskip
  \wosmlabel{Here: $\{p, q\}$ and $\{r\}$ by output, and nothing splits, so $p$ and
  $q$ merge. The identical-rows rule finds a subset of these merges.}
\end{frame}

\begin{frame}{Try it}
  \begin{enumerate}
    \item Draw a Moore machine whose output is 1 when its last two inputs were both
      1.\par
      \uncover<2->{\small Three states: no 1 yet (out 0), one 1 (out 0), two or more
      1s (out 1). Every 0 goes back to the first.}
    \item Now draw it as a Mealy machine. How many states does it need?\par
      \uncover<3->{\small Two: \enquote{last input 0} and \enquote{last input 1}. In
      the second, input 1 outputs 1.}
    \item Simplify this Moore machine: $p$: $0 \to q$, $1 \to r$, out 0; $q$:
      $0 \to p$, $1 \to s$, out 1; $r$: $0 \to q$, $1 \to r$, out 0; $s$: $0 \to p$,
      $1 \to s$, out 1.\par
      \uncover<4->{\small $p \equiv r$ and $q \equiv s$: two states are enough.}
  \end{enumerate}
\end{frame}

% =============================================================
\section{Summary}

\begin{frame}{Summary}
  \begin{itemize}
    \item Sequential circuits use the past as well as the present. The past is kept
      as a state.
    \item A finite state machine: finitely many states, an initial state, a
      transition rule and an output rule.
    \item Moore machines take the output from the state; Mealy machines from the
      state and the input.
    \item As a circuit: state memory, next-state logic and output logic.
    \item Equivalent states can be merged. Identical rows are a quick test;
      partition refinement finds every merge.
  \end{itemize}
\end{frame}

\begin{frame}{Before the next lecture}
  \begin{itemize}
    \item Read the beginning of Chapter~5 and Section~5.1 of the book.
    \item Draw the turnstile as a Mealy machine whose output says whether the arm
      turns.
    \item Simplify the book's Moore machine in Table 5.3 with partition refinement,
      and check that you get the same three states.
  \end{itemize}

  \begin{block}{Next lecture \unitref{11}}
    The clock, and memory: the SR latch and the D latch.
  \end{block}
\end{frame}

\end{document}
